\documentclass[11pt]{article}
\usepackage[margin=1in]{geometry}
\usepackage{amsmath,amssymb,amsthm}
\usepackage{hyperref}
\usepackage{xurl}
\let\oldus\_
\renewcommand{\_}{\oldus\allowbreak}

\newtheorem{theorem}{Theorem}
\newtheorem{lemma}[theorem]{Lemma}
\newtheorem{corollary}[theorem]{Corollary}
\theoremstyle{definition}
\newtheorem{definition}[theorem]{Definition}
\theoremstyle{remark}
\newtheorem{remark}[theorem]{Remark}

\newcommand{\F}{\mathbb{F}}

\title{A disproof of Conjecture 00000002913\\[2pt]
\large For every $d\ge2$ and every $q=2^k>d^2$, the Keller map $(x^2-x+y^d,\,y)$ of degree $d$ is not injective on $\F_q^2$}
\date{}

\begin{document}
\sloppy
\maketitle

\section{The conjecture}

Conjecture 00000002913 reads (English, verbatim):
\begin{quote}
Definition: finite field this of injection polynomial map: jacobian conjecture of finite field
plane. Conjecture: injection spectrum: d -fold polynomial map in F\_q on injection when q greater
than d$^2$ and inverse in q of injection does not exceed of d by Chebyshev lattice gives.
\end{quote}
The Chinese text is clearer. It defines the objects as injective polynomial maps of the
finite-field plane (``the finite-field version of the Jacobian conjecture'') and conjectures:
``a degree-$d$ polynomial map is injective on $\F_q$ when $q>d^2$, \emph{and} the failure of
the converse of injectivity for $q\le d$ is given by Chebyshev--Lagrange inversion.''

\paragraph{Which clause we refute.} The conjecture is a conjunction. We refute its first
conjunct, the \emph{threshold clause}: for all $d$ and $q$ with $q>d^2$, every degree-$d$
polynomial map over $\F_q$ is injective. We refute it under each of the following readings of
``polynomial map'':
\begin{itemize}
\item[(R1)] any polynomial map $\F_q^2\to\F_q^2$ of degree $d$ (no Jacobian hypothesis);
\item[(R2)] a \emph{Keller map} of the plane: a polynomial map $\F_q^2\to\F_q^2$ of degree $d$
whose Jacobian determinant is a nonzero constant (the hypothesis of the Jacobian conjecture, which
the definition names);
\item[(R3)] a one-variable polynomial map $\F_q\to\F_q$ of degree $d$.
\end{itemize}
Since (R2) is a special case of (R1), one family refutes both. The second conjunct (the ``$q\le d$''
clause) is not given a meaning and is not needed: a conjunction with a false conjunct is false.

\section{Definitions}

A polynomial map of the plane over a field $K$ is a pair $f=(f_0,f_1)$ of polynomials in
$K[x_0,x_1]$ (Lean: \texttt{PlaneMap K := Fin 2 -> MvPolynomial (Fin 2) K}). Its map on points is
$x\mapsto(f_0(x),f_1(x))$ (\texttt{PlaneMap.toFun}, evaluation by \texttt{MvPolynomial.eval}); its
degree is $\max(\deg f_0,\deg f_1)$ with $\deg$ the total degree (\texttt{PlaneMap.degree},
\texttt{MvPolynomial.totalDegree}); its Jacobian determinant is
$\det(\partial f_i/\partial x_j)$ with Mathlib's formal partial derivative
\texttt{MvPolynomial.pderiv} (\texttt{PlaneMap.jacobian}). Injectivity is
\texttt{Function.Injective}. $\F_q$ for $q=p^k$ is Mathlib's \texttt{GaloisField p k}, whose
cardinality is $p^k$ for $k\ne0$ (\texttt{GaloisField.card}).

\section{Proof}

For a prime $p$ and $m\ge1$ let
\[
  f_{p,m}(x,y) = (x^p - x + y^m,\; y)
\]
(Lean: \texttt{asMap p m}).

\begin{lemma}[\texttt{asMap\_not\_injective}]
Over every field, $f_{p,m}(0,0)=(0,0)=f_{p,m}(1,0)$; hence $f_{p,m}$ is not injective.
\end{lemma}
\begin{proof}
$0^p-0+0^m=0$ and $1^p-1+0^m=0$ since $p,m\ge1$.
\end{proof}

\begin{lemma}[\texttt{asMap\_jacobian}]
Over a field of characteristic $p$, the Jacobian determinant of $f_{p,m}$ is the constant $-1$.
\end{lemma}
\begin{proof}
$\partial_x(x^p-x+y^m)=p\,x^{p-1}-1=-1$, $\partial_y(y)=1$, $\partial_x(y)=0$, so the
determinant is $(-1)\cdot1-\partial_y(x^p-x+y^m)\cdot 0=-1$.
\end{proof}

\begin{lemma}[\texttt{totalDegree\_first}, \texttt{asMap\_degree}]
For $p\ge2$ and $m\ge1$, $\deg f_{p,m}=\max(p,m)$.
\end{lemma}
\begin{proof}
The upper bound follows from $\deg(a\pm b)\le\max(\deg a,\deg b)$. The monomials $x^p$ and $y^m$
both have coefficient $1$ (they differ from each other and from $x$), so the total degree is at
least $p$ and at least $m$. The second component $y$ has degree $1$.
\end{proof}

\begin{theorem}[(R1); \texttt{plain\_counterexample}]
For every field $K$ and every $d\ge2$, the map $(x^2-x+y^d,\,y)$ has degree exactly $d$ and is not
injective on $K^2$. In particular the threshold clause under (R1) fails for every $d\ge2$ and every
finite field $\F_q$ with $q>d^2$, prime fields included.
\end{theorem}

\begin{theorem}[(R2); \texttt{keller\_counterexample}]
Over every field of prime characteristic $p$ and for every $m\ge1$, $f_{p,m}$ has degree
$\max(p,m)$, Jacobian determinant $-1$, and is not injective.
\end{theorem}

\begin{theorem}[\texttt{conjecture\_00000002913\_false}]
Let $p$ be prime, $m\ge1$, and $k$ with $p^k>\max(p,m)^2$. Over $\F_q=$ \texttt{GaloisField p k}
one has $q=p^k$, $d:=\deg f_{p,m}=\max(p,m)$, $d^2<q$, Jacobian determinant $-1$, and $f_{p,m}$ is
not injective on $\F_q^2$.
\end{theorem}

\begin{corollary}[\texttt{keller\_counterexample\_every\_degree}, \texttt{sq\_lt\_two\_pow\_two\_mul}]
For every $d\ge2$ and every $k$ with $2^k>d^2$ (for example $k=2d$), the map $(x^2-x+y^d,\,y)$
over $\F_{2^k}$ has degree exactly $d$, Jacobian determinant $-1$, and is not injective. So for
every $d\ge2$ the threshold clause fails at infinitely many $q>d^2$, under (R1) and (R2), and also
for odd $d$ (where the characteristic does not divide the degree).
\end{corollary}

\begin{theorem}[(R3); \texttt{univariate\_counterexample}, \texttt{univariate\_derivative}]
For every field $K$ and $d\ge2$, $x^d-x\in K[x]$ has degree $d$ and $x\mapsto x^d-x$ is not
injective on $K$ ($0$ and $1$ are roots). If $\operatorname{char}K=p$ then $(x^p-x)'=-1$.
\end{theorem}

\begin{remark}[Context]
That $x\mapsto x-x^p$ defeats the naive Jacobian conjecture in characteristic $p$ is classical. The
Wikipedia article ``Jacobian conjecture'' states: ``For fields of positive characteristic, the
usual Jacobian conjecture is false in one dimension'' and, after giving the polynomial $x-x^p$,
``There are a few ways to modify the Jacobian conjecture to get around this issue.'' The
additional term $y^m$ only serves to realize every degree $d\ge2$ in the plane.
\end{remark}

\section{Scope and unrefuted readings}

Refuted: the threshold clause under (R1) and (R3) for every $d\ge2$ and every finite field with
$q>d^2$; under (R2) for every $d\ge2$ over $\F_{2^k}$ with $2^k>d^2$, and for every prime $p$ over
$\F_{p^k}$ with $p^k>\max(p,m)^2$. Not refuted:
(a) the \emph{existential} reading ``for $q>d^2$ some degree-$d$ map is injective'', which is true
(for instance the triangular automorphism $(x+y^d,y)$);
(b) reading (R2) restricted to \emph{prime} fields $\F_p$ (our Keller family needs $q=p^k$ with
$k\ge3$; under (R1) prime fields are covered by \texttt{plain\_counterexample});
(c) a reading that adds Adjamagbo's separability hypothesis ($p$ does not divide
$[K(x,y):K(f)]$), under which $f_{p,m}$ is excluded;
(d) $d=1$, where Keller maps are affine bijections and the clause is true;
(e) the second conjunct (the ``$q\le d$'' / Chebyshev--Lagrange inversion clause), which we do not
interpret.

\section{Relation to the earlier submission}

PR \#278 (orionsheep) used the reading (R3) with $x\mapsto x^2$ on $\F_5$ ($5>4$, $1^2=4^2$). It
was merged and then removed in the re-audit of commit 541cf4fb with the reason ``residue
numerals; no field/map/injectivity objects'': its Lean theorem was a conjunction of statements
about natural numbers such as $(4\cdot4)\bmod 5=(1\cdot1)\bmod5$. The present file uses Mathlib's
finite fields \texttt{GaloisField p k} with their cardinality, polynomial maps as
\texttt{MvPolynomial}/\texttt{Polynomial} evaluation, \texttt{Function.Injective}, total degree
and the formal Jacobian determinant. It refutes the clause for every degree $d\ge2$ and also under
the Jacobian (Keller) hypothesis that the definition names, which the earlier example could not
satisfy (the derivative of $x^2$ is not constant).

\section{Lean formalization}

File \texttt{lean/Conjecture2913/Basic.lean} (Lean 4.33.1, Mathlib), namespace
\texttt{Tlmc2913}. Main theorems: \texttt{conjecture\_00000002913\_false},
\texttt{keller\_counterexample\_every\_degree}, \texttt{keller\_counterexample},
\texttt{plain\_counterexample}, \texttt{univariate\_counterexample},
\texttt{univariate\_derivative}. \texttt{lake build} succeeds; \texttt{\#print axioms} of each gives
\texttt{[propext, Classical.choice, Quot.sound]}. No \texttt{sorry}, no \texttt{native\_decide}, no
added axioms.

\end{document}
